%%%%%%%%%%%%%%%%%%%%%%%%%%%%%%%%%%%%
%% Classe        %%%%%%%%%%%%%%%%%%%%%%%%%%%%
%%%%%%%%%%%%%%%%%%%%%%%%%%%%%%%%%%%%
\documentclass[12pt]{book}

%%%%%%%%%%%%%%%%%%%%%%%%%%%%%%%%%
%Quelques ajustements et variables   pour les marges%%%%
%%%%%%%%%%%%%%%%%%%%%%%%%%%%%%%%%
%L'origine par défaut correspond au coin en haut a gauche de la zone de texte
%Si on met \hoffset = 0 et \voffset = 0 : on sera a 1 inch du bord, et 1 inch du top c'est la valeur par défaut on va le changer pour partir du coin de la page.
\setlength{\hoffset}{-1in}% Offset horizontal global du texte pour être à 0
\setlength{\voffset}{-1in} % Offset vertical global du texte pour être à 0


%%%%%%%%%POUR LA HAUTEUR%%%%%%%%% A4= 29.7cm
\setlength{\topmargin}{1.5cm} % Marge en haut
\setlength{\headheight}{1.cm} % Taille du cadre pour le header
\setlength{\headsep}{.5cm}% Séparation entre le header et le texte( donné à la zone de texte header)
\setlength{\topskip}{0cm} % Séparation entre le header et le texte (donné à la 
%zone de texte) même effet que \headsep
\setlength{\textheight}{23cm} % Hauteur de la zone de texte (25cm)
\setlength{\footskip}{3cm} %Séparation avec la zone de footnote? % Bas de page + séparation %Marche pas

%%%%%%%%%POUR LA LARGEUR%%%%%%%%% A4= 21cm
%Si je met 3.25 pour odd et even, mon texte est centré (gauche droite)
%Je vais le décaler de +-0.25mm pour prendre en compte la reliure

\setlength{\oddsidemargin}{3.25cm}  % Marge gauche sur pages impaires
\setlength{\evensidemargin}{3.25cm} % Marge gauche sur pages paires

% largeur hauteur du texte
\setlength{\textwidth}{15cm} % Largeur de la zone de texte


%%%%%%%%%%%%%%%%%%%%%%%%%%%%%%%%%%%%
%% Packages   %%%%%%%%%%%%%%%%%%%%%%%%%%%%
%%%%%%%%%%%%%%%%%%%%%%%%%%%%%%%%%%%%
\usepackage[frenchb]{babel}
\usepackage[utf8]{inputenc}
\usepackage{pifont}
\usepackage[T1]{fontenc}
\usepackage{verbatim}
\usepackage[nooneline]{subfigure}
\usepackage{graphicx}% Include figure files
\usepackage{dcolumn}% Align table columns on decimal point
\usepackage{bm}% bold math
\usepackage{mathrsfs}
\usepackage{amsmath}
\usepackage{amssymb}
%\usepackage{geometry}
\usepackage{gensymb}
\usepackage{textcomp}
\usepackage{wasysym}
\usepackage{enumerate}
\usepackage{appendix}
\usepackage{lscape}

%Pour les petit table of content en début de chapitre
\usepackage{minitoc}
\mtcsettitle{minitoc}{Contenu}
%Profondeur d'indexage
\setcounter{secnumdepth}{4}
\setcounter{minitocdepth}{2}

\usepackage{caption}
\usepackage{amsmath}
\usepackage{lineno} % Numéro de lignes
\usepackage{ upgreek } %Lettres grec grandes
\usepackage{textcomp}
\usepackage[colorlinks= false, urlcolor=blue, citebordercolor={0 1 0}, linkbordercolor ={1 0 0},pdfborder={0 0 0}]{hyperref} %Gestion des affichages (Citations, Acronymes)
\usepackage[printonlyused,withpage]{acronym}  %Utilisation des acronymes

 \usepackage[version=3]{mhchem} %Pour les élements chimiques

\usepackage{multirow}
\usepackage{colortbl}
\usepackage{colortbl}
\usepackage[table]{xcolor}
\usepackage{array}
\usepackage{fancyhdr}
\usepackage{perpage}
\captionsetup{font=sl}


%%%%%%%%%%%%%%%%%%%%%%%%%%%%%%%%%%
%Commande pour modifier les haut de page %
%%%%%%%%%%%%%%%%%%%%%%%%%%%%%%%%%%
\usepackage{fancyhdr}



%%%%%%%%%%%%%%%%%%%%%%%%%%%%%%%%%%
%Commande pour avoir les footnote partant de 1 sur chaque page %
%%%%%%%%%%%%%%%%%%%%%%%%%%%%%%%%%%
\usepackage{perpage}
%\usepackage[bottom]{footmisc}% Pour aligner les footnote depuis le bas
\MakePerPage{footnote}


\usepackage{caption}%[2008/04/01]
\captionsetup{font=sl}

%%%%%%%%%%%%%%%%%%%%%%%%%%%%%%%%%
%Commande spéciale pour passer TableOfContent en Français %
%%%%%%%%%%%%%%%%%%%%%%%%%%%%%%%%%
\renewcommand{\contentsname}{Table des mati\`eres}

%%%%%%%%%%%%%%%%%%%%%%%%%%%%%%%%%
%Déclaration de toutes mes commandes   %%%%%%%%%%%%
%%%%%%%%%%%%%%%%%%%%%%%%%%%%%%%%%
\input Fonctions/myCommands.tex
\input Fonctions/ltxdefs.tex



%%%%%%%%%%%%%%%%%%%%%%%%%%%%%%%%%
% Page de couverture %%%%%%%%%%%%%%%%%%%%%%
%%%%%%%%%%%%%%%%%%%%%%%%%%%%%%%%%
\makeatletter

\def\clap#1{\hbox to 0pt{\hss #1\hss}}%

\def\ligne#1{
  \hbox to \hsize{
    \vbox{\centering #1}}}
    
\def\haut#1#2#3{
  \hbox to \hsize{
    \rlap{\vtop{\raggedright #1}}
    \hss
    \clap{\vtop{\centering #2}}
    \hss
    \llap{\vtop{\raggedleft #3}}}}
    
\def\bas#1#2#3{
  \hbox to \hsize{
    \rlap{\vbox{\raggedright #1}}
    \hss
    \clap{\vbox{\centering #2}}
    \hss
    \llap{\vbox{\raggedleft #3}}}}
    
% Titre
\def\maketitle{
  \thispagestyle{empty}\vbox to \vsize{
    \vfill
    \haut{}{\large \@blurb}{}
    \vfill
    \ligne{\LARGE \@title}
    \vspace{5mm}
    \ligne{\LARGE \@author}
    %\vspace{1mm}\ligne{\texttt{<\@email>}}
%    \vspace{1cm}
 %   \ligne{\Large \@administratif}
    \vfill
    \ligne{\Large \@jury}
    \vfill
%    \bas{

%   }
% % \cleardoublepage
	}
}

% Date, lieu,...
\def\date#1{\def\@date{#1}}
\def\author#1{\def\@author{#1}}
\def\title#1{\def\@title{#1}}
\def\location#1{\def\@location{#1}}
\def\blurb#1{\def\@blurb{#1}}
\def\jury#1{\def\@jury{#1}}
\def\email#1{\def\@email{#1}}
\def\administratif#1{\def\@administratif{#1}}
%\def\includegraphics#1{\def\@includegraphics{#1}}

% Valeurs par defaut
\date{\today}
\author{}
\title{}
\location{Villeurbanne}
\blurb{}

\email{steen@ipnl.in2p3.fr}

\makeatother
\title{\Huge Étude des gerbes hadroniques à l’aide du prototype du calorimètre hadronique semi-digitial et comparaison avec les modèles théoriques utilisés dans le logiciel GEANT4}
\author{Arnaud \textsc{Steen}}
\location{Villeurbanne}
\administratif{
  %Sous la direction de :\\
  %Imad  \textsc{Laktineh} et 
}
\blurb{
  Th\`ese de l'Universit\'e de Lyon \\
  \bigskip
  \'Ecole Doctorale de Physique et Astrophysique de Lyon (PHAST)\\[1em]
  \bigskip
  Th\`ese de Doctorat en Physique des Particules \\
  \bigskip
  Institut de Physique Nucl\'eaire de Lyon (IPNL)\\
}
\jury{
  \begin{center}
    \begin{tabular}{lll}
      JURY~: & M.  Didier \textsc{Contardo} & (examinateur)\\
      & M.  Gérald \textsc{Grenier} & (co-directeur de th\`ese)\\
      & M.  Imad  \textsc{Laktineh} & (directeur de th\`ese)\\
      & M.  Dominique \textsc{Pallin}       & (rapporteur)\\
      & M.  Alberto \textsc{Ribon}       & (examinateur)\\
      & M.  Laurent \textsc{Serin}      & (rapporteur)\\
      & M.  Frank \textsc{Simon}      & (examinateur)\\
    \end{tabular}
  \end{center}
}
%%%%%%%%%%%%%%%%%%%%%%%%%%%%%%%

%Pour que les footnotes n'aillent pas sur deux pages
\interfootnotelinepenalty=10000


%%###############################################
%%#######Début manuscrit     #########################
%%###############################################
\begin{document}

%%%###############################################
%%#######LOGOs   #################################
%%###############################################
\begin{figure}
  \begin{center}
    \includegraphics[width=0.18\textwidth]{Logo/logo_white_ILC.jpg}
    \hspace{1cm}
    \includegraphics[width=0.18\textwidth]{Logo/calice_logo.jpg}
    \hspace{1cm}
    \includegraphics[width=0.18\textwidth]{Logo/Logo_IPNL.jpg}
    \hspace{1cm}
    \includegraphics[width=0.18\textwidth]{Logo/logo_ucbl_lyon1.jpg}
  \end{center}
\end{figure}
	
\maketitle{\thispagestyle{empty}}

\clearpage



%%%%%%%%%%%%%%%%%%%%%%%%%%%%%%%%%
%Activation numéro de lignes            %%%%%%%%%%%%%%%
%%%%%%%%%%%%%%%%%%%%%%%%%%%%%%%%%
%\setpagewiselinenumbers
\modulolinenumbers[5]
\linenumbers


%%%%%%%%%%%%%%Haut et Bas de Page%%%%%%%%%%
\pagestyle{fancy}
\fancyhf{}
\renewcommand{\sectionmark}[1]{ \markright{\thesection\, #1}}
\renewcommand{\chaptermark}[1]{%
\markboth{\sc\thechapter.\ #1}{}}
%\renewcommand{\headrulewidth}{0.5pt}
%\renewcommand{\footrulewidth}{0.5pt}
%\renewcommand{\sectionmark}[1]{\markright{\thesection{} #1}}
\lhead[\thepage]{\rightmark}
 \rhead[\leftmark]{\thepage}
\fancyhead{} % clear all header fields
\fancyhead[RE,LO]{\thepage}
\fancyhead[RO,LE]{\rightmark}
\fancyfoot{\thepage} % clear all footer fields
\fancyfoot[LE,RO]{\thepage} 
\lfoot[\thepage]{} 
\rfoot[]{\thepage}


%%###############################################
%%#######Sauts de pages pour présentation  #############
%%###############################################
%\input Specials/PageVierge.tex
%\newpage
%\cleardoublepage

\dominitoc

%%%###############################################
%%%####### Resumé  ########################
%%%###############################################
%\frontmatter %active les numérot de pages en i
\pagestyle{empty}
\input Resume/Resume.tex
\pageblanche


%%%###############################################
%%%####### Remerciement  ########################
%%%###############################################
%%\frontmatter
%%\renewcommand*\thesection{\arabic{section}}
%%\pagestyle{empty}
%\thispagestyle{empty}
%\input Remerciement/Remerciement.tex

%\renewcommand*\thesection{\arabic{section}}
\input Introduction/Introduction.tex

%%###############################################
%%#######Table des Matières  ########################
%%###############################################

\tableofcontents{\thispagestyle{empty}}
\pageblanche
%\cleardoublepage

%%###############################################
%%#######Liste des Figures  ########################
%%###############################################
%\renewcommand\listfigurename{Liste des figures}
%\listoffigures{\thispagestyle{empty}}

%%###############################################
%%#######Liste des Tables  ########################
%%###############################################
%\listoftables{\thispagestyle{empty}}

%On désactive ici les haut de pages
%\pagestyle{empty}

%###############################################
%%#######MicroIntro###############
%%###############################################

%%%%%%%%%%%%%%Haut et Bas de Page%%%%%%%%%%
\pagestyle{fancy}
\fancyhf{}
%\renewcommand{\sectionmark}[1]{ \markright{\thesection\, #1}}
%\renewcommand{\chaptermark}[1]{\markboth{\ #1}{}}
%\renewcommand{\headrulewidth}{0.5pt}
%\renewcommand{\footrulewidth}{0.5pt}
%\renewcommand{\sectionmark}[1]{\markright{\thesection{} #1}}
\lhead[\thepage]{\rightmark}
 \rhead[\leftmark]{\thepage}
%\fancyhead{} % clear all hea

%###############################################
%%#######Chapitre: Introduction générale###############
%%###############################################

\mainmatter % désactive les numérot de pages en i

%On active ici les haut de pages comme définit dans fancy
\pagestyle{fancy}

\input ModelStandard/ModelStandard.tex
\input ILC/ILC.tex
\input ShowerTh/ShowerTh.tex
\input SDHCAL/SDHCAL.tex
\input Digitizer/Digitizer.tex
\input Shower/Shower.tex
\input ILD/ILD.tex

\backmatter
\input Conclusion/Conclusion.tex

%###############################################
%#######Bibliographie #############################
%###############################################
\bibliography{Bibliography/Bibliography}
\bibliographystyle{Bibliography/MyRefStyle}
\addcontentsline{toc}{chapter}{\bibname} 

  
  
\end{document}
